\documentclass[12pt]{article}
\usepackage[utf8]{inputenc}
\usepackage[T1]{fontenc}
\usepackage{hyperref}
\title{R\=a\d zi\d h\=i Language Resources for K-12 Curriculum Development}
\author{Compiled by Claude Sonnet 5.5 (claude-sonnet-5-5) for Andr\'as Kornai}
\date{}
\begin{document}
\maketitle

\section{Introduction}
R\=a\d zi\d h\=i (ISO 639-3 rzh, Glottocode jaba1234) is spoken around Jabal Razih in the Sa\textquoteleft dah area of northwestern Yemen. Sources give about 63,000 speakers, an Ethnologue figure from 2004. It may descend from Old South Arabian, or it may be a heavily influenced variety of Arabic, and sources disagree on this. I found no school materials, primers or community-language programmes.

\section{Primary Educational Resources}
None found. The language is not known to be taught in schools.

\section{Linguistic and Reference Materials}
\begin{itemize}
\item Wikipedia article on the Razihi language. \url{https://en.wikipedia.org/wiki/Razihi_language}
\item Omniglot page on Razihi. \url{https://omniglot.com/writing/razihi.htm}
\item Citation record for Watson, Glover Stalls, Razihi and Weir (2006), \emph{Two Texts from Jabal R\=a\d zi\d h, North-west Yemen}, the main scholarly text publication. \url{https://ociana.osu.edu/citations/11146}
\item Blog post by Charles H\"aberl on the vitality of Razihi. \url{https://blog.linguisten.de/post/642926024430059520}
\item Endangered Languages Project entry. \url{https://www.endangeredlanguages.com/lang/10839}
\end{itemize}

\section{Religious and Community Resources}
Joshua Project reports no Christian resources in Razihi. \url{https://joshuaproject.net/people_groups/20969}

\section{Digital and Audio Resources}
None found.

\section{Curriculum Development Context}
There is no known orthography, pedagogical grammar or reader. Jabal Razih is hard to reach, and the area has been affected by conflict, so fieldwork is difficult. Any curriculum would first need documentation and a community-agreed writing system. Shelagh Weir's \emph{A Tribal Order} (2007) was named in a search summary as a study of the speakers, but I saw no link for it.
\end{document}